\section{模型假设与符号说明}

\subsection{模型假设}
为保证模型可解并与题目信息一致，作如下假设：
\begin{enumerate}[label=\arabic*)]
    \item 所有车辆均从同一配送中心出发并最终返回配送中心；
    \item 路段速度由时段状态决定，可用分段期望速度近似；
    \item 客户服务时长固定（可取 $20$ 分钟）；
    \item 充电、加油等待时间在基础模型中忽略；
    \item 时间窗采用软约束，违反部分通过惩罚项计入目标函数。
\end{enumerate}

\subsection{符号定义}
\begin{table}[H]
\centering
\caption{核心符号说明}
\begin{tabular}{p{0.2\textwidth}p{0.65\textwidth}}
\toprule
符号 & 含义 \\
\midrule
$N=\{0,1,\dots,n\}$ & 节点集合，$0$ 为配送中心 \\
$K$ & 车辆集合 \\
$T$ & 车型集合 \\
$d_{ij}$ & 节点 $i$ 到节点 $j$ 的距离 \\
$q_i^w,q_i^v$ & 客户 $i$ 的重量、体积需求 \\
$[e_i,l_i]$ & 客户 $i$ 的服务时间窗 \\
$Q_t^w,Q_t^v$ & 车型 $t$ 的载重与容积上限 \\
$x_{ijk}^t$ & 是否由车型 $t$ 的车辆 $k$ 执行弧 $(i,j)$ \\
$y_{ik}^t$ & 客户 $i$ 是否由车型 $t$ 的车辆 $k$ 服务 \\
$t_i$ & 到达客户 $i$ 的时刻 \\
$z_k^t$ & 车型 $t$ 的车辆 $k$ 是否启用 \\
\bottomrule
\end{tabular}
\end{table}
